\documentclass[11pt,letterpaper]{article}
\usepackage[T1]{fontenc}
\usepackage[utf8]{inputenc}
\usepackage{lmodern,microtype}
\usepackage[margin=0.88in,headheight=15pt]{geometry}
\usepackage{amsmath,amssymb,amsthm,mathtools}
\usepackage{booktabs,longtable,tabularx,array}
\usepackage{xcolor,listings,xurl,hyperref,fancyhdr,enumitem,etoolbox}
\definecolor{ink}{HTML}{203B50}
\definecolor{papergray}{HTML}{F4F6F7}
\hypersetup{colorlinks=true,linkcolor=ink,citecolor=ink,urlcolor=ink,
 pdftitle={Asymptotic: A Differential Audit},
 pdfauthor={Technical review prepared for Vladimir Reshetnikov}}
\pagestyle{fancy}\fancyhf{}
\fancyhead[L]{\small\sffamily ASYMPTOTIC / DIFFERENTIAL AUDIT}
\fancyhead[R]{\small\sffamily 9 September 2026}
\fancyfoot[C]{\small\thepage}
\setlength{\parindent}{0pt}\setlength{\parskip}{5pt}
\setlength{\emergencystretch}{3em}
\renewcommand{\arraystretch}{1.15}
\newcolumntype{P}[1]{>{\raggedright\arraybackslash}p{#1}}
\newcommand{\code}[1]{\texttt{\detokenize{#1}}}
\newcommand{\OO}{\mathcal O}
\newcommand{\RR}{\mathbb R}
\newcommand{\NN}{\mathbb N}
\newcommand{\SHA}{921387e5ba1239bfda96e63e64e89bf63d9c41e6}
\newcommand{\repo}{https://github.com/VladimirReshetnikov/Asymptotic/blob/921387e5ba1239bfda96e63e64e89bf63d9c41e6/}
\newcommand{\source}[2]{\href{\repo#1}{\nolinkurl{#2}}}
\newtheorem{proposition}{Proposition}[section]
\newtheorem{lemma}[proposition]{Lemma}
\newtheorem{corollary}[proposition]{Corollary}
\lstset{basicstyle=\ttfamily\footnotesize,columns=fullflexible,
 keepspaces=true,showstringspaces=false,breaklines=true,breakatwhitespace=false,
 backgroundcolor=\color{papergray},frame=single,rulecolor=\color{black!15},
 framerule=0.3pt,xleftmargin=4pt,xrightmargin=4pt,aboveskip=8pt,belowskip=8pt,
 tabsize=2}
\BeforeBeginEnvironment{lstlisting}{\begin{samepage}}
\AfterEndEnvironment{lstlisting}{\end{samepage}}
\setcounter{tocdepth}{2}
\begin{document}
\begin{titlepage}
\thispagestyle{empty}
\vspace*{0.65in}
{\sffamily\large\color{ink} SOURCE REVIEW / NATIVE EXPERIMENTS / MATHEMATICAL ANALYSIS}\par
\vspace{0.3in}
{\sffamily\bfseries\Huge Asymptotic}\par
\vspace{0.14in}
{\sffamily\LARGE A Differential Audit}\par
\vspace{0.18in}
{\Large New evidence beyond the nine existing code reviews}\par
\vspace{0.4in}
{\large Wolfram Language comparison, a false-domain certificate,
flat-sector precision loss, and resolution-aware numerical checks}\par
\vspace{0.35in}
\begin{tabular}{@{}ll@{}}
Repository & \code{VladimirReshetnikov/Asymptotic}\\
Pinned revision & \texttt{\SHA}\\
Package version & 1.8.0\\
Native kernel used & Wolfram Language 15.0.0, Linux x86-64\\
Review date & 9 September 2026\\
Prepared for & Vladimir Reshetnikov
\end{tabular}
\vfill
\begin{minipage}{0.96\linewidth}
\textbf{Principal result.} A cleanly constructed inverse object for the source
domain $x>10$ can subsequently receive a successful certificate with root
$2$ and verification interval $[1,3]$ when unrelated ambient assumptions
are active. A narrowly scoped patch rejects this case and preserves a valid
control. Independently, flat-sector multiplication loses precision when one
input is made more informative; a proved, tested boundary-projection repair
restores the appropriate bound.

\textbf{Evidence boundary.} This report includes successful focused native
experiments, source analysis, and independently executed Python checks.
It does not claim a full-suite run, a complete formal verification, or measured
whole-package performance. Proposed tests and optimizations are explicitly
distinguished from executed checks.
\end{minipage}
\end{titlepage}

\begin{abstract}
This audit reviews a pinned revision of the Asymptotic repository while treating
its nine existing code reviews as an exclusion baseline. It contributes four
specific deltas. First, the previously recognized ambient-assumption problem
is sharpened to a native reproduction of an incorrect source-domain certificate
for an already constructed object. Second, the flat-sector product algorithm
prematurely erases exponential degrees: increasing one operand's sector depth
can worsen the algebraic factor in the reported tail. A boundary-sensitive
replacement is proved and tested. Third, a native numerical inverse check is
shown to return an inexact zero ratio with roughly minus 800 digits of accuracy;
this is not an exact-zero error, but it needs an explicit resolution status.
A normalized correction equation resolves the example without subtracting
nearly equal source roots. Fourth, native Wolfram 15 already supplies the first
three Dirichlet terms of the zeta function, contrary to an overbroad source
comment. The package's genuine additional value in that example is its explicit
coordinate and quantitative tail bounds. The report develops a concrete
fixed-order derivative bound for those tails, compares the broader architecture
with native capabilities, and supplies narrow patches, reproductions, test
specifications, and independent mathematical checks.
\end{abstract}
\tableofcontents
\newpage

\section{Executive assessment}
\label{sec:executive}
The repository is valuable as a real-asymptotic computation layer whose results
carry more structure than a printed approximation: selected branches,
local coordinates, retained coefficient blocks, remainder information,
source equations, and, on some paths, refinement state. Its scope is substantially
broader than its original inverse-power-series use case. The inspected sources
include ordinary power--log calculus, exact inverse cores, reciprocal-logarithmic
coordinates, Fourier-polynomial coefficients, finite flat sectors, native
special-function adapters, and exact rational interval certificates.
\cite{readme,seriesops,coordinates,fourier,flat,certificates}

It should not be evaluated against an artificially weak description of Wolfram
Language. Native \code{Series} already handles some fractional powers and
logarithms; \code{Asymptotic} uses scales beyond ordinary Taylor series;
\code{InverseSeries} performs reversion; and \code{AsymptoticSolve} supplies
local implicit solutions, including real-branch and multivariable problems.
The appropriate comparison is the complete contract needed by a workflow,
not the presence or absence of the word ``asymptotic.''\cite{wlseries,wlasymptotic,wlinverse,wlsolve}

\begin{longtable}{@{}P{0.55in}P{1.03in}P{3.08in}P{1.38in}@{}}
\toprule ID & Priority & New or materially sharper result & Evidence\\\midrule
\endfirsthead
\toprule ID & Priority & New or materially sharper result & Evidence\\\midrule
\endhead
N1 & High: certificate integrity & A stored domain is incorrectly marked verified
on an interval disjoint from it. The object was constructed outside the hostile
ambient context. & Native witness, clean control, patched rejection and valid control.\\
N2 & Medium: precision and wasted work & Flat-tail multiplication collapses degree
information too early. Depths $(1,1),(1,2),(1,3)$ produce tail powers $-4,-6,-8$
instead of the available $-1$. & Native baseline and patch; proof; 4,000 independent
bound-model checks.\\
N3 & Medium: diagnostic resolution & A numerical error ratio prints as zero but
has accuracy about $-799.89$. Direct subtraction cannot resolve the requested
scale at the working precision. & Native observation; independent 60- and
950-digit comparisons; stable correction-coordinate calculation.\\
N4 & Low: capability claims & Native \code{Asymptotic} and the finite part of
\code{Series} already recover the three-term zeta example. & Version-stamped
native comparison and package tail-bound comparison.\\
\bottomrule
\end{longtable}

These findings have different meanings. N1 is an incorrect certificate claim.
N2 returns a \emph{valid but unnecessarily weak} asymptotic bound, not a false
bound. N3 does not accuse Wolfram's arbitrary-precision arithmetic of treating
an inexact zero as an exact number: its accuracy metadata correctly reveals the
loss of information. N4 is a documentation and comparison correction, not a
reason to remove a useful dedicated adapter.

The highest-value development path from this audit is correspondingly narrow:
isolate the certificate proof context; retain exponential degree until the
last precision projection; make numerical resolution inspectable; and use
native algorithms as independent coefficient oracles on genuinely overlapping
families. The later derivative-tail construction is a concrete extension of
one existing family, not a repetition of a general transseries wishlist.

\section{Scope, reproducibility, and exclusion of old findings}
\label{sec:scope}
\subsection{What was examined and executed}
The review is pinned to \texttt{\SHA}. The standalone package downloaded by the
native evaluator contained 574,894 bytes and 38 inlined source modules. The
native runtime identified itself as
\begin{quote}\small
\code{15.0.0 for Linux x86 (64-bit) (May 6, 2026)}.
\end{quote}
The source investigation covered the review/status records, README, standalone
builder, series calculus and refinement paths, certificate and numerical-check
paths, source-coordinate reductions, flat and Fourier layers, reciprocal-log
operations, Gamma inverse operations, and several special-function layers.
The package was loaded using \code{URLDownload} followed by local \code{Get},
not direct evaluation of a compressed HTTP response.\cite{readme,builder,status}

The successful native observations were deliberately small. They include the
certificate witness and controls, flat products before and after the proposed
replacement, the numerical-resolution example, classical series/reversion
controls, zeta comparisons, and two direct calls to \code{AsymptoticSolve}.
The latter returned unevaluated on the tested irrational-power and flat-inverse
requests. That observation is not a proof that no native reformulation can solve
those problems.

Several other tool attempts returned HTTP 502 or network errors. Those attempts
provide no evidence that a package function is slow, incorrect, or unsupported.
No runtime or peak-memory number is inferred from a connector failure. The
accompanying JSON records successful observations as a faithful transcription,
not as a raw kernel transcript or a complete test report.

Independent local checks use SymPy for four flat-inverse coefficients and
mpmath for numerical comparisons. A small, separately implemented asymptotic
bound model tests 4,000 deterministic randomized projections and two exact-zero
edge cases. These checks validate the stated mathematical transformations and
model behavior; they do not execute the Wolfram package. An additional 28
high-precision numerical cases check the differentiated zeta-tail formula, and
17 synthetic-fixture checks exercise the patch utility.

\subsection{How duplication was controlled}
The maintained status ledger consolidates the existing reviews and identifies
both unresolved findings and recently completed fixes. It was used as the
primary exclusion map. Candidate topics were screened in the full article text
of reviews 1--6, 8, and 9; the relevant detailed findings and engineering
sections of review 7 were read. Zeta mentions were screened across all nine,
including review 7's mathematical-capabilities section. This was targeted
novelty screening, not a claim to have re-proved every earlier report.
\cite{reviewindex,status,r1,r7findings,r7engineering,r8}

\begin{longtable}{@{}P{0.64in}P{2.35in}P{3.05in}@{}}
\toprule Delta & Closest existing material & Additional contribution here\\\midrule
\endfirsthead
\toprule Delta & Closest existing material & Additional contribution here\\\midrule
\endhead
N1 & C05; review 1 A04 and review 8 F01: ambient assumptions can escape retained
provenance. & A clean object, changed ambient state only during verification,
and an explicitly false whole-interval domain assertion. A tested mitigation
at the proof step, with its limited scope stated.\\
N2 & Review 7's general information-monotonicity recommendation; P06's resource
budget discussion. & A precise flat-tail mechanism, an infinite family of
weakening bounds, actual outputs, a proof of the repair, and an implementation
replacement.\\
N3 & D05's checked-evaluation discussion; the existing separation between
numerical evidence and certification. & A concrete loss of about 800 digits
in the \emph{inverse-check} ratio, correct interpretation of the inexact zero,
and a stable normalized correction equation.\\
N4 & Review 1 lists zeta/Lerch defining-sum constructions. & An executed native
comparison showing the same finite terms, while distinguishing the package's
stronger tail metadata.\\
Extension & X04's broad differentiation/re-expansion opportunity. & An explicit
all-fixed-order derivative-tail formula for the Dirichlet family, with conditions,
proof, and a standalone Wolfram function.\\
\bottomrule
\end{longtable}

This report does \emph{not} re-report the dense native-export allocation problem,
the nested fractional-power guard, logarithmic-degree loss in native export,
generic fixed-margin refinement, a general resource-budget redesign, the head
rename/migration issue, CI proposals, license metadata, the primitive flat-rate
GCD proposal, or broad complex/multivariate/uniform-asymptotic extensions.
In particular, the ledger records completed fixes for the dense-export and
fractional-power issues; these are not presented as current discoveries.
\cite{status}

\subsection{Evidence vocabulary}
``Observed'' denotes a successful native or independent execution recorded in
the bundle. ``Source-derived'' denotes a consequence of inspected control flow.
``Proved'' denotes a mathematical argument written below. ``Proposed'' denotes
a change whose full integration remains work for the repository. A proposed
regression file is not silently counted as an executed test suite.

\section{Architecture and comparison with built-in Wolfram Language}
\label{sec:comparison}
\subsection{Compare mathematical contracts, not only finite expressions}
An asymptotic result can be regarded as a tuple
\[
 (\text{germ},\ \text{scale},\ \text{finite part},\ \text{error class},\
  \text{conditions},\ \text{evidence}).
\]
Two outputs can have the same finite expression and substantially different
usefulness. One may retain the original equation, another may be easier to
compose with native \code{SeriesData}, and a third may carry an explicit
pointwise tail bound. Conversely, a more elaborate result object is not better
when its additional asserted evidence is wrong. N1 illustrates the latter
case; N4 illustrates the former.

The ordinary series layer records a representation of the form
\[
 b(y)+a(y)\left(\sum_{\beta<H}w(y)^\beta P_\beta(\log w(y))
                 +\OO\!\left(w(y)^\rho M(w(y))^d\right)\right),
 \quad M(u)=1+|\log u|,
\]
with a positive local coordinate and an exact prefactor. The source explicitly
separates relative jet precision from that prefactor. Other families deliberately
use other structures: a retained Lambert core, finite Fourier modes, or separate
exponential-sector and inner-power truncations. They should not be flattened
into one native rational-lattice object merely because their finite parts can
be printed as a sum.\cite{seriesops,reciprocal,gammaops,flatops}

Native \code{SeriesData} represents powers on a rational lattice through an
integer denominator. That statement is about the data object's exponent
encoding, not a claim that \code{Series} can return only a single simple
\code{SeriesData} or cannot produce factored expressions. The observed native
zeta output is a useful counterexample to the latter misconception.
\cite{wlseriesdata}

\subsection{Capability map}
\begin{longtable}{@{}P{1.10in}P{2.31in}P{2.63in}@{}}
\toprule Task & Native capabilities & Repository contribution / boundary\\\midrule
\endfirsthead
\toprule Task & Native capabilities & Repository contribution / boundary\\\midrule
\endhead
Ordinary local series & \code{Series}, \code{SeriesCoefficient}, \code{Normal},
composition and reversion of suitable series. & Explicit power--log models and
retained source/remainder metadata; ordinary controls should agree.\\
Fractional powers and logs & \code{Series} handles some fractional and logarithmic
expansions; \code{Asymptotic} infers richer scales. & An explicit ordered
real-exponent algebra and complete logarithmic blocks, rather than requiring
one rational lattice for all exponents.\\
Implicit/inverse expansions & \code{InverseSeries}; local solutions from
\code{AsymptoticSolve}, with real/complex and multivariable forms. & Specialized
real inverse methods, branch/source information, power observables, and
refinement of supported nonclassical scales.\\
Exact inverse cores & Native \code{ProductLog} and special functions are useful
building blocks. & Retains an exact core while ordering corrections relative to
it; avoids replacing the core by a numerically inadequate finite log expansion.\\
Gamma/Barnes inverse families & Native forward asymptotics and Lambert functions
can support hand-derived inversion. & Family-specific inversion and coefficient
blocks; supported calculus is intentionally narrower than ordinary jets.\\
Oscillatory log coefficients & Native trigonometric/exponential algebra and
asymptotic tools can aid derivations. & A finite Fourier-polynomial coefficient
algebra with exact frequencies and explicit frequency budgets.\\
Flat sectors & Native expressions can represent exponentials; individual
asymptotic and transformed problems may be solvable. & A separate sector degree,
exact zero sector, inner bounds, and flat-specific operations; N2 concerns
precision inside this existing machinery.\\
Special functions & Broad native symbolic function coverage and asymptotics. &
A real-domain adapter layer and selected explicit identities/tail models;
function coverage alone does not establish a unique advantage.\\
Zeta defining sum & Native 15.0.0 supplies the tested three finite terms through
both \code{Asymptotic} and \code{Normal[Series[...]]}. & Explicit $w=e^{-x}$ scale,
positive lower tail, and a quantitative upper tail bound.\\
Numerical comparison & Arbitrary precision, \code{FindRoot}, and precision/accuracy
tracking. & Inverse-aware checks and retained original equations; N3 needs an
additional resolution classification, not a replacement for native arithmetic.\\
Interval evidence & Native exact arithmetic and symbolic processing are usable
primitives. & A custom outward-rational enclosure implementation and a local
root certificate; its domain proof must be isolated from ambient state.\\
General analysis & Native asymptotic integral, transform, equation and related
interfaces cover a broad surrounding system. & A specialized library, not a
replacement for that surrounding system. No new blanket adapter wishlist is
proposed here.\\
\bottomrule
\end{longtable}
\noindent Sources for this map are the public guide and inspected modules,
together with official documentation for the native functions.
\cite{readme,seriesops,coordinates,reciprocal,gammaops,fourier,flat,flatops,special,
identities,dirichlet,numerical,certificates,wlseries,wlinverse,wlsolve,wlasymptotic}

\subsection{Native controls that prevent unfair comparisons}
Two inexpensive native controls were executed:
\begin{lstlisting}
Series[x^x, {x, 0, 3}]
InverseSeries[Series[x + x^2, {x, 0, 5}], y]
\end{lstlisting}
Their returned coefficient data are respectively
\begin{align*}
 &\{1,\log x,(\log x)^2/2,(\log x)^3/6\},\qquad O(x^4)\text{ endpoint},\\
 &\{1,-1,2,-5,14\},\qquad y-y^2+2y^3-5y^4+14y^5+O(y^6).
\end{align*}
The first is a native \code{SeriesData} whose displayed coefficients themselves
contain logarithms. It directly rules out the claim that logarithmic
coefficients alone make the package unique. It is not used here to infer a
stronger mathematical logarithmic-tail contract than the native object states.

A more specialized test was
\begin{lstlisting}
AsymptoticSolve[x + x^Sqrt[2] == y,
  {x, 0}, {y, 0, 3}, Reals]
\end{lstlisting}
which remained unevaluated on the tested kernel. The package's explicit
inverse constructor with the exclusive cutoff 3 returned
\begin{align*}
 g(y)={}&y-y^{\sqrt2}+\sqrt2\,y^{2\sqrt2-1}
 +\left(-3+\frac1{\sqrt2}\right)y^{3\sqrt2-2}\\
 &+\left(-4+\frac{17\sqrt2}{3}\right)y^{4\sqrt2-3}
 +\text{its recorded remainder}.
\end{align*}
Every displayed exponent is below 3. The direct native request for
$x+e^{-1/x}=y$ at the positive zero endpoint likewise remained unevaluated,
whereas the dedicated flat constructor supplied sector expansions used below.
These are useful concrete capability deltas, but not claims of impossibility
for native transformations, manually introduced markers, or other algorithms.
The numeral 3 is not a matched-order benchmark: the package call uses an
exclusive exponent cutoff, whereas the native implicit solver's approximation
order is not necessarily a polynomial-degree cutoff. The comparison records
what the stated direct requests do, not equivalence of their order conventions.

\subsection{Sound design choices worth preserving}
The source contains substantial safeguards that should survive optimization.
A generic magnitude remainder is not automatically differentiable. The ordinary
calculus checks derivative contracts. Exact source transformations are treated
differently from approximate exponentiation. Gamma and Barnes inverse scales
have explicit operation restrictions instead of pretending to be ordinary
polynomial-log jets. The Fourier layer checks real conjugacy relations and
uses frequency limits as hard resource bounds rather than silently deleting
modes.\cite{seriesops,coordinates,gammaops,fourier}

The exact special-function normalization layer deserves particular credit.
For example, a half-integer modified Bessel function can contain both dominant
and subdominant exponentials even when a dominant asymptotic series terminates.
The implementation normalizes exact identities before finite asymptotic
truncation, and it retains the second exponential where justified. Similar
care appears in reflected error functions and finite hypergeometric identities.
A review should not propose ``recognize exact termination'' without preserving
that distinction.\cite{identities}

The standalone builder also has meaningful engineering safeguards: it parses
relevant source-level imports, records module hashes, rejects repeated/cyclic
or out-of-scope dependencies, and validates the assembled result before an
atomic replacement. These features are present, not new recommendations.
\cite{builder}

\section{N1: a false source-domain certificate from ambient assumptions}
\label{sec:n1}
\subsection{The witness}
Construct the inverse object under a neutral ambient context:
\begin{lstlisting}
s = Block[{$Assumptions = True},
  AsymptoticInverse`AsymptoticInverse[
    ConditionalExpression[x, x > 10],
    {x, Infinity}, {y, 3}]];
\end{lstlisting}
The object retains \code{SourceDomain -> x > 10}. Now perform only the certificate
call inside a different context:
\begin{lstlisting}
Block[{$Assumptions = x > 10},
  AsymptoticInverse`InverseCertificate[s, 2,
    "Interval" -> {1, 3}, "Center" -> 2,
    "MaxRefinements" -> 0]]
\end{lstlisting}
The unpatched native result included
\begin{lstlisting}
"Certified" -> True
"RootEnclosure" -> {2, 2}
"VerificationInterval" -> {1, 3}
"CertifiedSourceDomain" -> x > 10
"SourceDomainVerified" -> True
"CertifiedErrorBound" -> 0
\end{lstlisting}
The same call with \code{$Assumptions = True} returned
\code{Failure["OutsideBranch", ...]}. Thus the defect is not explained by an
already corrupted object or by an ambiguous source branch chosen at
construction.

The distinction between the equation and its domain is crucial. The number
2 really is a root of the unrestricted equation $x=2$, so its residual and
zero error bound for that unrestricted equation are unsurprising. But there
is no root of the restricted problem $x=2$, $x>10$. Moreover, \emph{every} point
of the verification interval violates the retained domain. The incorrect claim
is the certificate's branch/domain assertion, not the elementary arithmetic.

\subsection{Mechanism and logical error}
The relevant function is \code{certPositiveCondition} in
\path{InverseCertificates.wl}. Before interval comparisons, it performs
\begin{lstlisting}
c = TimeConstrained[
  Quiet[Refine[condition, Element[x, Reals]]],
  1, condition];
\end{lstlisting}
and accepts \code{c === True}. Passing a second assumption argument to
\code{Refine} does not, by itself, neutralize ambient assumptions. The following
native probe returned \code{True}; analogous probes of \code{Simplify} and
\code{FullSimplify} did too:
\begin{lstlisting}
Block[{$Assumptions = x > 10},
  Refine[x > 10, Element[x, Reals]]]
\end{lstlisting}
This is normal assumption-aware simplification behavior. The mistake is using
its conclusion as a statement about all points of an independently specified
interval.\cite{certificates,wlrefine,r1,r8}

The proof obligation is
\[
 \forall t\in[1,3],\quad t>10.
\]
The simplifier is instead allowed to use a hypothesis equivalent to the
proposition being checked. The valid implication
\[
 x>10\ \Longrightarrow\ x>10
\]
does not establish the quantified interval claim. The error occurs before the
otherwise useful exact-rational interval arithmetic becomes relevant.

This is why the new witness is materially stronger than merely showing that a
constructor records \code{Assumptions -> True} after using a parameter
hypothesis. The proof stage itself can turn a false branch assertion into a
successful evidence record for an unchanged, correctly constructed object.
Adding the hostile assumption to the stored domain would not fix this
quantifier error: the domain would still be false on the interval.

\subsection{Narrow mitigation, observed controls, and its limits}
The supplied patch changes the simplification step to
\begin{lstlisting}
c = Block[{$Assumptions = True},
  TimeConstrained[
    Quiet[Refine[condition, Element[x, Reals]]],
    1, condition]];
\end{lstlisting}
The exact old anchor occurred once in the downloaded standalone file. A
semantically identical replacement was applied in the native test. The hostile
$[1,3]$ call then returned \code{Failure["OutsideBranch", ...]} with
\code{Certified -> False}. A positive control with target 12, center 12, and
interval $[11,13]$ still returned a certificate with enclosure $\{12,12\}$.

The patch is intentionally described as a \emph{mitigation of this proof-step
leak}, not a repository-wide cure for ambient assumptions. The broad
construction/operation-context problem is already recorded in C05. The supplied
change neither captures all public-entry hypotheses nor validates every other
symbolic oracle in the package.

For the certificate subsystem, the stronger integration invariant should be:
\begin{quote}
After an inverse object and exact numerical target are fixed, changing ambient
assumptions must not create additional unconditional certificate claims.
Only explicitly admitted and correctly scoped evidence may discharge a proof
obligation.
\end{quote}
In particular, a predicate being verified on an interval must not be available
as an unexamined pointwise hypothesis for its own proof. Parameter conditions
must either be checked after specialization or remain explicit conditions of a
conditional result; they cannot be silently treated as numerical interval facts.

\subsection{Tests that specifically exercise the proof boundary}
The bundle supplies the false-domain witness, a clean-context control, and a
valid interval control. Further integration tests should vary the domain's
Boolean structure: strict and nonstrict inequalities, conjunctions,
disjunctions, negations, rational endpoints, and predicates whose truth differs
between interval endpoints. These are not a second generic request for
assumption tests: their acceptance criterion is the truth of a whole-interval
certificate after construction has already completed.

A particularly useful negative control uses ambient assumptions that conflict
with the interval rather than with the stored source expression. It catches
precisely the difference between symbolic implication and universal interval
containment. Tests should inspect the structured certificate fields, not only
whether an expression simplifies or an exception is raised.

\section{N2: exponential-degree erasure weakens flat products}
\label{sec:n2}
\subsection{Representation and the current projection}
The flat layer uses an exponential marker
\[
 E(u)=\exp(-c/u^p),\qquad c,p>0,\quad u\to0^+,
\]
and separate finite power--log coefficients. An object with sector depth $N$
retains sectors $0,\ldots,N$ and a complete tail of the form
\[
 E(u)^{N+1}\OO\!\left(u^\rho M(u)^d\right).
\]
Unknown inner coefficient errors remain in their original sectors. This is
an important distinction already implemented by the package.\cite{flat,flatops}

For input depths $N_a,N_b$, \code{flatOpsMultiplyData} chooses
$N=\min(N_a,N_b)$, computes the finite sector convolution, and collects all
omitted contributions into a tail beginning at $E^{N+1}$. Its current loop
replaces every extra exponential factor by the inequality $E^q\le1$ and then
combines the remaining algebraic bounds. The same loss happens for a tail
times a nonzero sector and for a product of two tails.\cite{flatops}

This is conservative but can be extremely wasteful. A factor $E^q$ with
$q>0$ is not merely bounded by 1: it beats \emph{every fixed} inverse power
and logarithmic factor. Erasing $q$ before comparing algebraic exponents loses
exactly the information needed to avoid polluting the first omitted sector.

\subsection{A reproduced anti-refinement family}
Let $g$ be the positive local inverse of
\[
 f(x)=x+e^{-1/x},\qquad x\to0^+,
\]
and put $E=e^{-1/y}$. Formal inversion, or direct substitution, gives
\begin{align}
 g(y)={}&y-E+y^{-2}E^2
 +\frac{y-3/2}{y^4}E^3
 +\frac{y^2-4y+8/3}{y^6}E^4+\cdots.
 \label{eq:flat-inverse}
\end{align}
The four coefficients were independently generated by
\[
 C_k(y)=\frac{(-1)^k}{k!}e^{k/y}
             \frac{d^{k-1}}{dy^{k-1}}e^{-k/y}.
\]
For the argument below, only the first two coefficients are necessary.

Construct $s_1$ with depth 1 and $s_N$ with depth $N$, then multiply them:
\begin{lstlisting}
a = AsymptoticInverse`AsymptoticFlatInverse[
  x + Exp[-1/x], {x, 0}, {y, 1}];
b = AsymptoticInverse`AsymptoticFlatInverse[
  x + Exp[-1/x], {x, 0}, {y, N}];
p = AsymptoticInverse`FlatSeriesMultiply[a, b];
\end{lstlisting}
The native outputs for $N=1,2,3$ all have the same finite expression,
\[
 y^2-2yE,
\]
but their tail pairs are
\[
 \{-4,0\},\qquad\{-6,0\},\qquad\{-8,0\},
\]
respectively, each multiplied by $E^2$. Increasing the information in the
second operand makes the reported common-depth product \emph{less} precise.

For this model the constructor's tail at depth $N$ is
\[
 E^{N+1}\OO(y^{-2N}).
\]
The product of the two input tails therefore contributes
\[
 E^{N+3}\OO(y^{-2N-2}).
\]
The current projection weakens this to $E^2\OO(y^{-2N-2})$, which dominates the
other projected algebraic envelopes and explains the observed progression.
Higher finite coefficients have lowest powers no worse than $-2N+2$; the
$k-1$ differentiations in $C_k$ lower the power by at most $2(k-1)$. Thus the
same mechanism predicts the family $\rho_{\rm old}=-2N-2$ for all $N\ge1$.

The available bound is much better:
\begin{equation}
 g(y)^2=y^2-2yE+E^2\OO(y^{-1}).
 \label{eq:sharp-product}
\end{equation}
Indeed, squaring~\eqref{eq:flat-inverse} shows that the coefficient of $E^2$ is
$1+2/y$. Equally, the depth-1 uncertainty $E^2\OO(y^{-2})$ multiplied by the
exact zero sector $y$ produces $E^2\OO(y^{-1})$. Terms with extra exponentials
do not require a worse algebraic power. The old and improved scale estimates
at $y=10^{-3}$, $N=3$, differ by the factor $10^{21}$. This is a comparison
of asymptotic scales, not of certified numerical constants.

\subsection{A boundary-sensitive projection theorem}
\begin{lemma}\label{lem:absorb}
For $\delta,c,p>0$ and fixed real $a,b$,
\[
 e^{-\delta c/u^p}u^aM(u)^b\longrightarrow0
 \quad\text{as }u\to0^+.
\]
\end{lemma}
\begin{proof}
The logarithm of the positive expression is
$-\delta c/u^p+a\log u+b\log M(u)$. The first term tends to minus infinity
faster than the absolute values of the other two grow. Exponentiating gives
the claim.
\end{proof}

Write the two input objects abstractly as
\[
 A=\sum_{i=0}^{N_a}E^i A_i+E^{N_a+1}R_a,
 \qquad
 B=\sum_{j=0}^{N_b}E^j B_j+E^{N_b+1}R_b.
\]
Each coefficient may itself carry an inner power--log uncertainty. Bounds
below apply to the whole coefficient, not just its displayed finite part.
Set $N=\min(N_a,N_b)$ and $D_k=\sum_{i+j=k}A_iB_j$.

\begin{proposition}\label{prop:projection}
At common retained depth $N$, the algebraic envelope at the first omitted
sector $E^{N+1}$ need only combine the following contributions:
\[
 D_{N+1},\qquad
 \mathbf 1_{N_a=N}\,R_aB_0,\qquad
 \mathbf 1_{N_b=N}\,R_bA_0.
 \tag{*}
\]
If these have a finite power--log bound $\OO(u^\rho M^d)$, the complete omitted
product tail is $E^{N+1}\OO(u^\rho M^d)$. If every contribution in $(*)$ is
exactly zero but a later omitted contribution may be nonzero, the conservative
fallback $E^{N+1}\OO(1)$ is valid. An exact-zero tail is justified only when all
omitted contributions are exact zero.
\end{proposition}
\begin{proof}
Expand the product. Finite convolution terms omitted from the retained result
have degree at least $N+1$. A term $R_aB_j$ has degree $N_a+1+j$; it has degree
exactly $N+1$ precisely when $N_a=N$ and $j=0$. The analogous assertion holds
for $R_bA_i$. The tail product has degree $N_a+N_b+2>N+1$.

Every remaining contribution has at least one additional positive integer
power of $E$. By Lemma~\ref{lem:absorb}, its fixed power--log envelope can be
absorbed into any chosen finite power--log envelope at degree $N+1$. There
are finitely many coefficient contributions and finitely many tail-product
bounds. Their constants can therefore be combined. If the boundary is empty,
choose the envelope 1 rather than claiming that no later contribution exists.
\end{proof}

The proposition does not move an unknown inner error to a later sector. Such
an error is included in the bound of its original $A_i$ or $B_j$ and therefore
in the appropriate $D_k$. This is essential: simply forgetting inner errors
would be an unsound and quite different change.

\subsection{Patch behavior and independent validation}
The supplied replacement implements $(*)$ after the existing convolution.
It preserves the rest of the function, including compatibility checks,
coefficient arithmetic, domain combination, and derivative provenance. It
also explicitly handles an empty boundary with a possible later tail.

On the same three native cases, the patched routine returned the unchanged
finite expression and tail pairs
\[
 \{-1,0\},\qquad\{-1,0\},\qquad\{-1,0\}.
\]
The independent Python model separately enumerates all tail contributions with
their exponential degrees and compares them with the optimized boundary rule.
It passed 4,000 seeded random cases using rational power exponents, logarithmic
degrees, unequal depths, and exact zeros, plus two explicit empty-boundary
cases. This tests the projection algebra, not all of the upstream coefficient
and domain machinery.

The change is intentionally modest. It \emph{does not yet avoid} computing the
full finite convolution and therefore does not claim a runtime speedup.
Its purpose is to repair precision loss with a small reviewable change. The
next section separates the additional work-elimination opportunity.

\section{N3: a zero numerical ratio with no useful resolution}
\label{sec:n3}
\subsection{Observed output and its correct interpretation}
For the depth-1 inverse above, the native call
\begin{lstlisting}
AsymptoticInverse`InverseNumericalCheck[
  a, 1/1000, WorkingPrecision -> 50]
\end{lstlisting}
returned the following important fields:
\begin{lstlisting}
"Error" -> 0``62.69897000433602
"Ratio" -> 0``-799.8899938021677
"RootResidual" -> 0``63.
"RemainderScale" -> 2.57653587296114965... * 10^-863
"Approximation" -> 0.001`50.
"ReferenceRoot" -> 0.001`50.
\end{lstlisting}
Double-backtick notation specifies accuracy. The ratio is \emph{not} an exact
zero and is not a well-resolved zero. Its large negative accuracy means the
calculation has no useful information at the scale needed to compare that
ratio with a number of order 1. Native precision tracking is exposing, rather
than concealing, the cancellation.\cite{wlaccuracy}

The source uses \code{WorkingPrecision + 10} for the approximation and reference
root, then subtracts them. The remainder scale is evaluated afterward. It is
not used to decide whether the root subtraction has enough resolution. The
output's statement that it is numerical evidence rather than an interval
certificate is correct, but it does not replace a resolution assessment for
this particular diagnostic.\cite{numerical}

\subsection{The true error is strictly positive}
Let $y>0$ be small enough that $q=y-e^{-1/y}>0$, and set $E=e^{-1/y}$. Since
$q<y$,
\[
 f(q)-y=-E+e^{-1/q}<0.
\]
Also $f'(x)=1+e^{-1/x}/x^2>0$ for $x>0$, and $f(y)>y$. The unique positive
root of $f(x)=y$ is therefore strictly between $q$ and $y$. The true error
$g(y)-q$ is positive, independently of any floating-point computation.
Equation~\eqref{eq:flat-inverse} gives its leading scale:
\[
 g(y)-q=\frac{E^2}{y^2}\left(1+\OO(Ey^{-2})\right).
\]
At $y=1/1000$,
\[
 \log_{10}(E^2/y^2)=-862.5889638065036553\ldots.
\]
Subtracting two roots of magnitude $10^{-3}$ to obtain an error of magnitude
about $10^{-863}$ loses roughly 860 significant digits. A 50-digit requested
root computation with ten guard digits cannot resolve it.

The independent 60-digit direct subtraction returned zero. The 950-digit
calculation returned
\[
 2.576535872961149652190150499507352912722130573760632202117945887577139
 \times10^{-863},
\]
with error divided by $E^2/y^2$ equal to 1 at the displayed 70-digit precision.
These are numerical checks, not interval enclosures. Their interpretation is
supported by the separate strict-positivity argument.

\subsection{A resolution-aware public diagnostic}
An improved check should expose at least the accuracy of the subtraction and
of the normalized ratio, and distinguish a successful root solve from a
resolved asymptotic-error comparison. A result such as
\begin{lstlisting}
"ReferenceSolveStatus" -> "Converged"
"ErrorResolutionStatus" -> "Insufficient"
"RatioAccuracy" -> -799.8899...
"ResolvedErrorRatio" -> Missing["InsufficientResolution"]
\end{lstlisting}
would accurately summarize the native information already present.
This is a proposed diagnostic schema, not an existing API.

A simple adaptive fallback can estimate the extra working precision needed
from the observed ratio accuracy or, before subtraction, from
\[
 \log_{10}\!\frac{|g_{\rm approx}|}{\text{remainder scale}}
       +\text{desired ratio digits}+\text{guard margin}.
\]
For this example, roughly 870 digits are needed for a ten-digit ratio before
adding a guard margin. This is a workload-dependent planning estimate, not a
universal precision theorem. Repeated computations must confirm stabilization,
and a maximum precision budget must be respected. A formal $O$ scale is not
a pointwise certified upper bound, so the adaptation must not label the result
``certified.''

The better solution on some families is to avoid subtracting the source roots
at all. That opportunity is concrete here.

\subsection{Normalize the correction before solving}
Write
\[
 g(y)=y-E+\frac{E^2}{y^2}w,
 \qquad u=y-E+\frac{E^2}{y^2}w.
\]
The desired error ratio is exactly $w$. Define
\[
 t=-\frac{E(1-Ew/y^2)}{uy},\qquad
 \operatorname{exprel}(t)=\frac{e^t-1}{t},
 \quad \operatorname{exprel}(0)=1.
\]
Then $f(u)=y$ is equivalent, where the displayed denominators are nonzero, to
\begin{equation}
 w=\frac{y(1-Ew/y^2)}{u}\operatorname{exprel}(t).
 \label{eq:normalized}
\end{equation}
To verify it, observe that $-1/u+1/y=t$ and divide
$-E+E^2w/y^2+Ee^t=0$ by $E^2/y^2$. The factor $(e^t-1)/t$ removes the
cancellation analytically before numerical evaluation.

At $y=1/1000$, $w$ is of order 1 even though the source error is of order
$10^{-863}$. Using a stable \code{expm1} evaluation in the included Python
implementation, the normalized equation gives ratio 1 to the displayed
55 digits at 60-digit working precision. The direct 950-digit calculation
agrees. This is a change of numerical coordinates, not artificial padding of
input precision.

The method suggests a useful diagnostic extension: when an inverse object
retains a leading omitted sector or a controlled residual model, solve for a
normalized correction variable rather than always for the original source
root. The eligibility checks must establish a meaningful nonzero scale and
an appropriate local branch. This report supplies the exact scalar example
and its numerical implementation, not a claim of a universal automatic
normalization algorithm.

\section{N4: a version-stamped correction to the zeta comparison}
\label{sec:n4}
The introductory comment in \path{DirichletSpecialFunctions.wl} describes
expansions from defining sums as unavailable from native \code{Series}.
That is too broad for the inspected runtime. The following native request
returned $1+2^{-x}+3^{-x}$:
\begin{lstlisting}
Asymptotic[Zeta[x], x -> Infinity, SeriesTermGoal -> 3]
\end{lstlisting}
The native \code{Series[Zeta[x], {x, Infinity, 3}]} result had a less convenient
nested form,
\begin{lstlisting}
1 + 2^SeriesData[x, Infinity, {-1}, -1, 4, 1]
  + 3^SeriesData[x, Infinity, {-1}, -1, 4, 1]
\end{lstlisting}
but applying \code{Normal} again gave $1+2^{-x}+3^{-x}$.
The package's corresponding three-term request returned the same finite part.
These observations concern this particular version and input; they do not
claim a native result for every transformed zeta or Lerch problem.
\cite{dirichlet,wlasymptotic,wlseries}

The package additionally returned the positive local coordinate $w=e^{-x}$,
a lower remainder bound $4^{-x}$, and an upper bound
\begin{equation}
 4^{-x}\left(1+\frac4{x-1}\right),\qquad x>1.
 \label{eq:zeta-bound}
\end{equation}
These are meaningful additional outputs. They follow directly from the positive
defining tail and the integral test; they are not obtained merely by dropping
a nested native remainder.

The appropriate correction is therefore not ``delete the adapter because the
kernel can do it.'' Instead, change the capability wording, add a native
coefficient-equivalence test for this overlapping case, and document the
dedicated adapter's explicit error and coordinate contract. A benchmark should
separately compare finite-term generation and generation of the richer bound
record. Reporting one as a speed comparison with the other would compare
different tasks.

There is also an oracle-quality benefit. A dedicated defining-sum implementation
and the native \code{Asymptotic} route are independently useful for checking
finite coefficients. But agreement of their finite parts does not prove the
package's remainder bound or source-domain claim. Those require the integral
argument and the branch checks respectively.

\section{Performance opportunities derived from N2 and N3}
\label{sec:performance}
\subsection{Avoid computing coefficients that only need an envelope}
The current flat product computes all finite sector-pair products even though
its output retains only through $N=\min(N_a,N_b)$. Proposition~\ref{prop:projection}
shows that an improved implementation needs exact convolution coefficients
only through sector $N+1$: retained sectors and the first omitted boundary.
Higher sectors need only the fact that their coefficients have the supported
fixed power--log envelopes.\cite{flatops}

For the reproduced depth pair $(1,N)$, the current rectangular convolution has
$2(N+1)$ candidate pairs. For $N\ge2$, only five pairs have total degree at most
2. Three contribute to the retained sectors 0 and 1, and two contribute to
the boundary sector 2. This constant-size coefficient workload is independent
of how far the second operand was refined. The supplied precision patch still
uses the original full convolution; the five-pair algorithm is the proposed
second-stage optimization.

The difference matters when coefficients are large symbolic power--log
polynomials. Multiplying and expanding a deep coefficient merely to learn a
bound that an extra exponential will absorb is unnecessary work. It can also
cause a local resource refusal unrelated to the information the requested
output actually needs. This is a concrete demand-pruning opportunity, rather
than a repetition of the existing proposal to rename or split \code{MaxTerms}.

A safe implementation should retain the following facts without materializing
all later products: whether a coefficient is exact zero; its supported envelope
class; whether a tail may be nonzero; and the exponential degree at which a
contribution starts. It should not infer exact zero merely because the first
omitted coefficient happens to cancel. The empty-boundary branch in the patch
is specifically designed to avoid that error.

\subsection{Separate proof cost, coefficient cost, and diagnostic cost}
For this delta, three measurements would be informative: exact pair products
actually computed, coefficient leaf counts before and after expansion, and
the achieved tail pair. Wall time alone cannot detect a precision regression.
The baseline returns a weaker tail when the second input is deeper, so a
comparison that ignores the tail can falsely describe the computation as
successful refinement.

For numerical diagnostics, measure the precision needed for a resolved
\emph{ratio}, not just for a converged source root. The normalized equation
provides an example where the source-root coordinate is catastrophically
inefficient for the diagnostic but a correction coordinate is well scaled.
The independent 60/950-digit experiments establish the distinction in required
precision; this audit does not claim a measured speedup factor.

The most useful acceptance comparison is therefore multi-dimensional:
identical branch and finite part, an equal or stronger remainder, a resolved
diagnostic when claimed, and reduced work. Existing broad performance and
benchmarking recommendations need not be repeated to implement these specific
comparisons.

\section{A concrete extension: differentiable Dirichlet tail contracts}
\label{sec:dirichlet-extension}
The existing roadmap already mentions differentiation and re-expansion.
Here is a specific addition with a closed-form bound, rather than another
request for a general differentiation engine. It applies to the existing
zeta defining-sum family.\cite{dirichlet,status}

For integer $m\ge2$ and $s>1$, let
\[
 R_m(s)=\sum_{n=m}^{\infty}n^{-s}.
\]
On every compact subset of $s>1$, fixed-order differentiated series converge
uniformly, since $(\log n)^j$ grows more slowly than every positive power of
$n$. Termwise differentiation gives
\[
 (-1)^j R_m^{(j)}(s)
   =\sum_{n=m}^{\infty}(\log n)^j n^{-s}=:T_{m,j}(s)>0.
\]
The nonnegativity is useful evidence in its own right.

\begin{proposition}\label{prop:zeta-derivative}
For an integer $j\ge0$, $m\ge2$, $s>1$, and $s\log m\ge j$,
\begin{align}
 T_{m,j}(s)\le{}&m^{-s}(\log m)^j\\
 &+m^{1-s}\sum_{k=0}^{j}
       \frac{j!}{(j-k)!}
       \frac{(\log m)^{j-k}}{(s-1)^{k+1}}.
 \label{eq:zeta-derivative-bound}
\end{align}
The right side is an explicit upper bound, not merely an asymptotic scale.
\end{proposition}
\begin{proof}
For $h(t)=t^{-s}(\log t)^j$, the logarithmic derivative is
\[
 \frac{h'(t)}{h(t)}=\frac1t\left(\frac{j}{\log t}-s\right).
\]
Thus $h$ is nonincreasing for $t\ge m$ under the stated condition. The integral
test bounds its sum by $h(m)+\int_m^{\infty}h(t)\,dt$. With $v=\log t$,
\[
 \int_m^{\infty}t^{-s}(\log t)^j\,dt
 =\int_{\log m}^{\infty}e^{-(s-1)v}v^j\,dv.
\]
Repeated integration by parts gives the finite sum in
\eqref{eq:zeta-derivative-bound}. Equivalently, the integral is
$\Gamma(j+1,(s-1)\log m)/(s-1)^{j+1}$, with the upper incomplete gamma
function. The finite expression avoids introducing another special-function
bound evaluator.
\end{proof}

For $j=0$, the formula is exactly
\[
 R_m(s)\le m^{-s}\left(1+\frac{m}{s-1}\right),
\]
which specializes to the observed package bound~\eqref{eq:zeta-bound} at $m=4$.
For an affine source $s=ax+b$, the absolute $j$th derivative with respect to
$x$ acquires the factor $|a|^j$. The real source must still lie in the domain
where the sum and monotonicity bound apply.

\subsection{Do not confuse source derivatives with local-coordinate derivatives}
The package's ordinary derivative calculus is expressed in its recorded local
coordinate. If $w=e^{-s}$, then
\[
 \frac{d^j}{dw^j}w^{\log n}
 = (\log n)^{\underline j}w^{\log n-j},
\]
where the underline denotes a falling factorial. It is not correct to attach
a local-coordinate derivative contract merely because an $x$-derivative bound
was written down.

For every fixed $j$, the absolute falling factorial is bounded by a polynomial
of degree $j$ in $\log n$ with fixed nonnegative coefficients. Applying
Proposition~\ref{prop:zeta-derivative} term by term to that polynomial gives,
eventually as $w\to0^+$,
\[
 \frac{d^j}{dw^j}R_m(-\log w)
       =\OO\!\left(w^{\log m-j}\right).
\]
This is the form compatible with the ordinary power-coordinate derivative
contract. It supports every \emph{fixed} derivative order; no uniform assertion
for $j\to\infty$ is being made. For an affine $s(x)$, multiplication by the
coordinate Jacobian supplies the expected source-variable scale.

The supplied \path{zeta_tail_derivatives.wl} implements the explicit
source-exponent bound with its conditions. It treats $j=0$ separately and,
for positive $j$, starts at $\max(m,2)$ so the $n=1$ derivative vanishes without
creating a $0^0$ corner case. It is an independent proposal, not a patched
public package API or an interval certificate. Integrating the contract should
be restricted to the proved Dirichlet family rather than indiscriminately
raising derivative orders on all special-function adapters.

\section{Remediation plan and integration acceptance criteria}
\label{sec:plan}
\subsection{First: protect certificate semantics}
Apply and review the narrow N1 change, regenerate the standalone artifact, and
run the focused false-domain and valid-domain tests on both modular and
standalone loading paths. The broader C05 work remains tracked in its existing
location. The new acceptance target is that a neutral object cannot gain a
false unconditional domain certificate merely by changing ambient state during
verification.

A full release gate remains the repository's existing maintenance topic; this
report does not claim its focused controls replace one. The narrow patch should
not be marketed as fixing every hypothesis-capture issue.

\subsection{Second: repair the flat tail before optimizing its computation}
Use the boundary projection to preserve the strongest directly available
first-omitted-sector envelope. Validate unequal depths, nontrivial zero sectors,
inner remainders, cancellations at the first omitted coefficient, exact scalar
zero, logarithmic factors, and negative powers in the supported common chart.
The proof covers fixed supported envelopes and commensurate sector degrees;
it does not add new phase families.

Only after this precision behavior is stable should exact coefficient work
above the first omitted boundary be pruned. The precise acceptance condition
is unchanged retained coefficients and equal or stronger justified tails,
with fewer unnecessary products. Request-budget behavior must remain explicit
when a needed coefficient is genuinely expensive.

\subsection{Third: classify diagnostic resolution}
Preserve the existing numerical/certified distinction and additionally label
whether the reported ratio is numerically resolved. A converged reference root
and an unresolved error comparison should coexist in the result rather than
one being mistaken for the other. Supply an adaptive-precision route with a
hard limit and a correction-coordinate route only where a justified
normalization is available.

The normalized flat example is a suitable regression because both the true
error sign and its first scale are analytically known. A test must not merely
assert that the printed ratio is zero or that \code{FindRoot} returned a
number. It should examine numerical accuracy and agreement with the
independent normalized coordinate.

\subsection{Fourth: correct and strengthen native comparisons}
Update the zeta comment and add the version-stamped overlapping-family check.
Keep raw native \code{Series} output as well as its finite part: silently
normalizing away all native remainder information would make a misleading
representation comparison. For the package, retain its explicit upper/lower
bounds and their conditions in the record.

The separate irrational-power and flat-inverse probes are useful positive
examples of the package's practical scope. They should be recorded as direct
native calls that remained unevaluated on a specific version, not as universal
proofs that the native system cannot solve any equivalent formulation.

\subsection{Then: add the concrete derivative-tail contract}
Review Proposition~\ref{prop:zeta-derivative}, test the finite bound expression
against high-precision differentiated defining tails in its admitted domain,
and connect its local-coordinate form to the existing derivative-contract
mechanism. Keep the argument and admissibility conditions alongside the
implementation. This is an achievable extension of a present module, not a
new general transseries platform.

\section{Artifacts and how to reproduce the work}
\label{sec:artifacts}
The archive contains the article source and compiled PDF, two exact-source
patch replacements, a local-checkout patch utility, a native reproduction
script, corrected-behavior test specifications, the derivative-bound proposal,
an independent Python check program, and evidence/novelty records.

The patch utility defaults to a dry run. It checks the audited Git revision
unless explicitly overridden, requires each source anchor exactly once, refuses
to overwrite its backup, and writes each changed file atomically. It changes
only a local checkout; it does not push commits or edit a remote repository.
Run from the extracted bundle, for example:
\begin{lstlisting}
python code/apply_patches.py /path/to/Asymptotic
python code/apply_patches.py /path/to/Asymptotic --write
# In that checkout, regenerate with the repository's documented
# validation/build_standalone.py workflow and check freshness.

wolframscript -file code/reproduce.wl \
  /path/to/Asymptotic/AsymptoticInverse.wl

python code/independent_checks.py
\end{lstlisting}
The native script records the loaded package path, hash, and kernel version.
Run it once against the pinned baseline and once against the regenerated
patched build. The complete \path{acceptance.wlt} file is a supplied integration
specification; its individual core witnesses were exercised here, but the file
as a whole was not run as a test suite during this audit.

The JSON evidence file is explicitly a transcription of successful connector
outputs. The independent-results file is generated by the included Python
program. Their separation prevents a mathematical oracle check from being
mistaken for native execution of the repository. The novelty matrix identifies
the closest old finding for each new contribution and records what was excluded.

\section{Conclusion}
The repository's strongest direction is not to compete with every built-in
symbolic-analysis feature. It is to provide reliable reusable results in selected
real asymptotic scales, with explicit conditions and honest evidence. That
position is well supported by the inspected architecture and by the practical
irrational-power and flat-inverse examples.

The new work needed here is specific. A certificate must not borrow the very
predicate it is meant to prove from ambient state. A flat-tail product must not
erase exponential degree before deciding the algebraic envelope. A numerical
comparison must not leave users to interpret an apparently zero ratio whose
accuracy says it is unresolved. And native capability claims should be checked
against current native outputs, not inferred from the inconvenience of their
representation.

The supplied patches and proofs address two of these boundaries directly.
The normalized correction example and the differentiated Dirichlet-tail bound
show how further development can produce concrete mathematical improvements
without repeating the existing broad roadmap. Their common principle is to
preserve the information that gives an asymptotic result its meaning until the
last operation that actually needs to discard it.

\appendix
\section{Compact evidence ledger}
\begin{longtable}{@{}P{1.08in}P{3.48in}P{1.10in}@{}}
\toprule Check & Result & Status\\\midrule
\endfirsthead
\toprule Check & Result & Status\\\midrule
\endhead
Ambient oracle & \code{Refine}, \code{Simplify}, \code{FullSimplify} each simplify
$x>10$ to True under ambient $x>10$. & Native observed.\\
N1 baseline & Restricted inverse, target 2, interval $[1,3]$: certificate succeeds
only under the hostile ambient context. & Native observed.\\
N1 patch & Hostile invalid interval rejected; target 12 in $[11,13]$ still certified. & Native observed.\\
N2 baseline & Depth pairs $(1,1),(1,2),(1,3)$ yield $E^2$ tail powers $-4,-6,-8$. & Native observed.\\
N2 patch & Same finite parts, tail power $-1$ in all three cases. & Native observed.\\
Bound model & 4,000 seeded random comparisons plus two exact-zero edge cases. & Independent, executed.\\
N3 baseline & Inexact zero ratio with accuracy $-799.8899938021677$. & Native observed.\\
N3 normalization & 60-digit normalized ratio agrees with direct 950-digit
subtraction at displayed precision. & Independent, executed.\\
Native controls & Logarithmic \code{Series}; classical \code{InverseSeries};
three zeta terms. & Native observed.\\
Specialized inverse & Two direct \code{AsymptoticSolve} calls remain unevaluated;
package constructs the tested results. & Native observed.\\
Patch release & Full cross-family regression, all loaders, full native suite. & Not performed here.\\
Derivative bound & Closed finite upper bound and proof with explicit conditions;
28 independent numerical spot checks. & Proposed integration; independent checks executed.\\
\bottomrule
\end{longtable}

\begin{thebibliography}{99}\small
\bibitem{readme} Asymptotic repository, \source{README.md}{README.md}, pinned revision \texttt{\SHA}.
\bibitem{status} Asymptotic repository, \source{docs/development/CODE_REVIEW_STATUS.md}{docs/development/CODE_REVIEW_STATUS.md}, consolidated status and cross-review mapping.
\bibitem{reviewindex} Asymptotic repository, \source{code-review/README.md}{code-review/README.md}, index of nine earlier reviews.
\bibitem{r1} Review 1, \source{code-review/code-review-1/article/article-standalone.tex}{code-review-1/article/article-standalone.tex}, especially A04.
\bibitem{r7findings} Review 7, \source{code-review/code-review-7/article/03-findings.tex}{code-review-7/article/03-findings.tex}, correctness boundaries and information monotonicity.
\bibitem{r7engineering} Review 7, \source{code-review/code-review-7/article/04-engineering-and-roadmap.tex}{code-review-7/article/04-engineering-and-roadmap.tex}.
\bibitem{r8} Review 8, \source{code-review/code-review-8/article/asymptotic_repository_audit.tex}{code-review-8/article/asymptotic_repository_audit.tex}, especially F01.
\bibitem{builder} Asymptotic repository, \source{validation/build_standalone.py}{validation/build_standalone.py}.
\bibitem{seriesops} Asymptotic repository, \source{AsymptoticInverse/Kernel/SeriesOperations.wl}{Kernel/SeriesOperations.wl}, especially \code{seriesData}, \code{seriesMake}, and \code{seriesDerivative}.
\bibitem{coordinates} Asymptotic repository, \source{AsymptoticInverse/Kernel/SourceCoordinates.wl}{Kernel/SourceCoordinates.wl}.
\bibitem{reciprocal} Asymptotic repository, \source{AsymptoticInverse/Kernel/ReciprocalLogOperations.wl}{Kernel/ReciprocalLogOperations.wl}.
\bibitem{gammaops} Asymptotic repository, \source{AsymptoticInverse/Kernel/GammaInverseOperations.wl}{Kernel/GammaInverseOperations.wl}.
\bibitem{fourier} Asymptotic repository, \source{AsymptoticInverse/Kernel/FourierCoefficients.wl}{Kernel/FourierCoefficients.wl}.
\bibitem{flat} Asymptotic repository, \source{AsymptoticInverse/Kernel/FlatSectors.wl}{Kernel/FlatSectors.wl}, especially \code{flatConstruct} and its tail formula.
\bibitem{flatops} Asymptotic repository, \source{AsymptoticInverse/Kernel/FlatSectorOperations.wl}{Kernel/FlatSectorOperations.wl}, especially \code{flatOpsMultiplyData} and \code{flatOpsJetBound}.
\bibitem{certificates} Asymptotic repository, \source{AsymptoticInverse/Kernel/InverseCertificates.wl}{Kernel/InverseCertificates.wl}, especially \code{certPositiveCondition} and the closed-interval verification path.
\bibitem{numerical} Asymptotic repository, \source{AsymptoticInverse/Kernel/NumericalInverseChecks.wl}{Kernel/NumericalInverseChecks.wl}, especially \code{numericalInverseEvidence}.
\bibitem{special} Asymptotic repository, \source{AsymptoticInverse/Kernel/SpecialFunctionAdapters.wl}{Kernel/SpecialFunctionAdapters.wl}.
\bibitem{identities} Asymptotic repository, \source{AsymptoticInverse/Kernel/SpecialFunctionIdentities.wl}{Kernel/SpecialFunctionIdentities.wl}.
\bibitem{dirichlet} Asymptotic repository, \source{AsymptoticInverse/Kernel/DirichletSpecialFunctions.wl}{Kernel/DirichletSpecialFunctions.wl}.
\bibitem{wlseries} Wolfram Research, \href{https://reference.wolfram.com/language/ref/Series.html}{\emph{Series}}, official documentation, accessed 9 September 2026.
\bibitem{wlseriesdata} Wolfram Research, \href{https://reference.wolfram.com/language/ref/SeriesData.html}{\emph{SeriesData}}, official documentation, accessed 9 September 2026.
\bibitem{wlinverse} Wolfram Research, \href{https://reference.wolfram.com/language/ref/InverseSeries.html}{\emph{InverseSeries}}, official documentation, accessed 9 September 2026.
\bibitem{wlasymptotic} Wolfram Research, \href{https://reference.wolfram.com/language/ref/Asymptotic.html}{\emph{Asymptotic}}, official documentation, accessed 9 September 2026.
\bibitem{wlsolve} Wolfram Research, \href{https://reference.wolfram.com/language/ref/AsymptoticSolve.html}{\emph{AsymptoticSolve}}, official documentation, accessed 9 September 2026.
\bibitem{wlrefine} Wolfram Research, \href{https://reference.wolfram.com/language/ref/Refine.html}{\emph{Refine}}, official documentation, accessed 9 September 2026; ambient behavior also tested directly.
\bibitem{wlaccuracy} Wolfram Research, \href{https://reference.wolfram.com/language/ref/Accuracy.html}{\emph{Accuracy}}, official documentation, accessed 9 September 2026.
\end{thebibliography}
\end{document}
